\section{Generalizations}

The CLT can be generalized in many ways. None of the
assumptions (independence, identical distributions,
even finite variance) are entirely necessary. A
central paradigm of probability theory is that any
random quantity that arises as a sum of many small
contributions that are either independent or not too
strongly dependent, will converge to the normal
distribution in some asymptotic limit. Thousands of
examples exist, but there is no single
all-encompassing theorem that includes all of them as
a special case. Rather, probabilists have a toolbox of
tricks and techniques that they try to apply in order
to prove normal convergence in any given situation.
characteristic functions are among the more useful
techniques. Another important technique, the so-called
$\textbf{moment method}$, involves the direct use of
moments: If we can show that
$\mathbf{E}(W_n^k)\to \mathbf{E}(Z^k)$, where
$(W_n)_{n=1}^\infty$ is the (normalized) sequence
being studied, and $Z\sim N(0,1)$, then by Theorem
3.3.12 in [Dur2010] (p. 105) or Theorem (3.12) in
[Dur2004] (p. 109), that implies that
$W_n\implies N(0,1)$.

We now discuss several examples of interesting
generalizations of CLT.

\subsection{Triangular arrays}

\begin{thm}[Lindeberg-Feller CLT for triangular arrays]
\label{thm-lindeberg-feller}
Let $(X_{n,k})_{1\le k\le n<\infty}$ be a triangular
array of r.v.'s. Denote $S_n = \sum_{k=1}^n X_{n,k}$
(the sum of the $n$-th row). Assume that:

\begin{enumerate}[label=\arabic*.]
\item For each $n$, the r.v.'s $(X_{n,k})_{k=1}^n$ are
independent.

\item $\mathbf{E} X_{n,k}=0$ for all $n,k$.

\item $\mathbf{V}(S_n) = \sigma_n^2 \to \sigma^2<\infty$ as
$n\to\infty$.

\item For all $\epsilon>0$,
${\displaystyle \lim_{n\to\infty}} \sum_{k=1}^n
\mathbf{E}\left( X_{n,k}^2
\mathbf{1}_{\{|X_{n,k}|>\epsilon\}}\right)=0$.
\end{enumerate}

\noindent Then $S_n \implies N(0,\sigma^2)$ as $n\to\infty$.
\end{thm}

\begin{proof}
See [Dur2010], p. 110--111 or [Dur2004], p. 115--116.
The proof uses the characteristic function technique
and is a straightforward extension of the proof for
the i.i.d. case.
\end{proof}

\begin{example}
(Record times and cycles in permutations). Let
$X_1, X_2, \ldots$ be i.i.d. $U(0,1)$ r.v.'s. Let
$A_n$ be the event that
$\{X_n = \max(X_1,\ldots,X_n)\}$ (in this case, we
say that $n$ is a $\textbf{record time}$). Let
$S_n = \sum_{k=1}^n \mathbf{1}_{A_k}$ be the number
of record times up to time $n$. We saw in a homework
exercise that the $A_k$'s are independent events and
$\mathrm{P}(A_k) = 1/k$. This implies that
$\mathbf{E}(S_n) = \sum_{k=1}^n \frac{1}{k} = H_n$
(the $n$-th $\textbf{harmonic number}$) and
$\mathbf{V}(S_n) = \sum_{k=1}^n \frac{k-1}{k^2}$.
Note that both $\mathbf{E}(S_n)$ and $\mathbf{V}(S_n)$
are approximately equal to $\log n$, with an error
term that is $O(1)$. Now taking
$X_{n,k} =
(\mathbf{1}_{A_k}-k^{-1})/\sqrt{\mathbf{V}(S_n)}$
in \hyperref[thm-lindeberg-feller]{Theorem~\ref*{thm-lindeberg-feller}}, it is easy to check
that the assumptions of the theorem hold. It follows
that

\[
\frac{S_n-H_n}{\sigma(S_n)} \implies N(0,1).
\]

\noindent Equivalently, because of the asymptotic behavior of
$\mathbf{E}(S_n)$ and $\mathbf{V}(S_n)$ it is also
true that

\[
\frac{S_n-\log n}{\sqrt{\log n}} \implies N(0,1).
\]

\noindent $\textbf{Note:}$$S_n$ describes the distribution of
another interesting statistic on random permutations.
It is not too difficult to show by induction (using an
amusing construction often referred to as the
$\textbf{Chinese restaurant process}$) that if
$\sigma \in S_n$ is a uniformly random permutation on
$n$ elements, then the number of cycles in $\sigma$
is a random variable which is equal in distribution to
$S_n$.
\end{example}

\subsection{Erdös-Kac theorem}

\begin{thm}[Erdös-Kac theorem (1940)]
Let $g(m)$ denote the number of prime divisors of an
integer $k$ (for example, $g(28)=2)$. For each
$n\ge 1$, let $X_n$ be a uniformly random integer
chosen in $\{1,2,\ldots,n\}$, and let $Y_n=g(X_n)$ be
the number of prime divisors of $X_n$. Then we have

\[
\frac{Y_n - \log\log n}{\sqrt{\log\log n}} \implies N(0,1).
\]

\noindent In other words, for any $x\in\R$ we have

\[
\frac{1}{n}\# \left\{ 1\le k\le n: g(k) \le \log\log n + k \sqrt{\log\log n} \right\} \xrightarrow[n\to\infty]{} \Phi(x).
\]

\end{thm}

\begin{proof}
See [Dur2010], p. 114--117 or [Dur2004], p. 119--124.
The proof uses the moment method.
\end{proof}

\noindent Note that $Y_n$ can be written in the form
$\sum_{p\le n} \mathbf{1}_{\{p | X_n\}}$, namely the
sum over all primes $p\le n$ of the indicator of the
event that $X_n$ is divisible by $p$. The probability
that $X_n$ is divisible by $p$ is roughly $1/p$, at
least if $p$ is significantly smaller than $n$.
Therefore we can expect $Y_n$  to be on the average
around

\[
\sum_{\textrm{prime }p\le n} \frac{1}{p},
\]

\noindent a sum that is known (thanks to Euler) to behave
roughly like $\log\log n$. The Erdös-Kac theorem is
intuitively related to the observation that these
indicators $\mathbf{1}_{\{p | X_n\}}$ for different
$p$'s are also close to being independent (a fact
which follows from the Chinese remainder theorem). Of
course, they are only approximately independent, and
making these observations precise is the challenge to
proving the theorem. In fact, many famous open
problems in number theory (even the Riemann
Hypothesis, widely considered to be the most important
open problem in mathematics) can be formulated in
terms of a statement about approximate independence
(in some loose sense) of some arithmetic sequence
relating to the prime numbers.

\subsection{The Euclidean algorithm}

As a final example from number theory, consider the
following problem: For some $n\ge 1$, choose $X_n$
and $Y_n$ independently and uniformly at random in
$\{1,2,\ldots,n\}$, and compute their greatest common
divisor (g.c.d.) using the Euclidean algorithm. Let
$N_n$ be the number of division (with remainder)
steps that were required. For example, if $X = 58$
and $Y=24$ then the application of the Euclidean
algorithm would result in the sequence of steps

\[
(58,24) \to (24,10) \to (10,4) \to (4,2) \to (2,0),
\]

\noindent so 4 division operations were required (and the
g.c.d. is 2).

\begin{thm}[CLT for the number of steps in the Euclidean algorithm; D. Hensley (1992)]
There exists a constant $\sigma_\infty$ (which has a
very complicated definition) such that

\[
\frac{N_n - \frac{12 \log 2}{\pi^2} \log n}{\sigma_\infty \sqrt{\log n}} \implies N(0,1).
\]

\end{thm}

\noindent Hensley's theorem was in recent years significantly
generalized and the techniques extended by Brigitte
Vallée, a French mathematician. The fact that the
average value of $N_n$ is approximately
$(12\log 2/\pi^2)\log n$ was previously known from
work of Heilbronn and Dixon in 1969--1970, using ideas
dating back to Gauss, who discovered the probability
distribution now called the ``Gauss measure''. This is
the probability distribution on $(0,1)$ with density
$\frac{1}{\log 2(1+x)}$, which Gauss found (but did
not prove!) describes the limiting distribution of the
ratio of a pair of independent $U(0,1)$ random
variables after many iterations of the
division-with-remainder step in the Euclidean
algorithm.
